\title{Direct Preference Optimization: Your Language Model is Secretly a Reward Model}

\begin{abstract}
Large-scale unsupervised language models (LMs) acquire extensive general knowledge and some degree of reasoning capability; however, exerting precise influence over their outputs remains challenging, largely due to their fully unsupervised training processes. Current approaches aiming for greater steerability rely on collecting human judgments comparing the relative merits of different outputs, subsequently fine-tuning the base LM to better reflect these human preferences—commonly through reinforcement learning from human feedback (RLHF). Yet, RLHF is frequently intricate and prone to instability, as it necessitates first training a reward model to capture human preferences, then leveraging reinforcement learning to adapt the base LM so that it maximizes this reward signal without straying significantly from its original state. In this work, we propose a novel reward model parameterization within the RLHF framework, which allows us to analytically derive the corresponding optimal policy in closed form. This approach enables us to resolve the standard RLHF objective through a straightforward classification loss. Our proposed method, termed \textit{Direct Preference Optimization} (DPO), is robust, effective, and efficient—removing the requirement to sample from the LM during fine-tuning or to carry out extensive hyperparameter searches. Experimental results demonstrate that DPO is capable of fine-tuning LMs to better accord with human preferences, often outperforming existing techniques. Remarkably, DPO-based fine-tuning surpasses PPO-based RLHF in sentiment control tasks, and achieves comparable or superior results in summarization and single-turn dialogue tasks, all while being significantly easier to implement and train.
\end{abstract}

\section{Introduction}
Extensive unsupervised training of large language models (LMs) on massive datasets leads to the emergence of unexpected capabilities~\citep{chowdhery2022palm, brown2020language, touvron2023llama,bubeck2023sparks}. However, these models are exposed to human-generated data encompassing a broad spectrum of intentions, priorities, and proficiency levels. Not all of these aspects are desirable for the model to replicate. For instance, while an AI assistant for programming may benefit from being able to \textit{recognize} frequent coding mistakes to provide corrections, it is preferable for its outputs to reflect the higher-quality—if less common—coding skills present in its training corpus. Similarly, a language model might need to be \textit{cognizant} of widespread misconceptions (e.g., those held by 50\% of individuals), but it should not propagate these errors by presenting them as correct in half of its responses regarding those topics. Thus, it is essential to selectively encourage \emph{targeted behaviors and responses} from the model's broad spectrum of \textit{acquired knowledge and abilities}, a requirement fundamental for constructing AI systems that are controllable, effective, and safe~\citep{ouyang2022training}. Prevailing strategies typically utilize reinforcement learning (RL) to align LMs with human preferences; in this work, we demonstrate that the RL-based preference objective currently employed can, in fact, be solved exactly with a straightforward binary cross-entropy formulation, thereby streamlining the preference alignment process.

Broadly speaking, current techniques shape the behaviors of language models through carefully curated datasets that reflect human judgments on what constitutes safe and useful behavior. This alignment via preference learning is conducted following a preliminary stage of unsupervised large-scale pretraining on textual data. While supervised fine-tuning using human-provided exemplary responses offers the most direct route for preference learning, the prevailing and most effective approach is reinforcement learning from human or AI feedback (RLHF/RLAIF; \citep{christiano2017deep,bai2022constitutional}). RLHF entails training a reward model based on datasets of human preferences and subsequently applying RL to optimize the LM so that its policy yields outputs rated highly by the reward model, all while constraining it from diverging significantly from the pretrained base model. Although RLHF leads to language models with advanced conversational and programming skills, the process is notably more involved than traditional supervised fine-tuning, requiring the simultaneous training of several language models and iterative policy sampling within the training loop, which substantially increases computational overhead.

This work presents a method for optimizing language models directly with respect to human preferences, sidestepping both explicit reward modeling and reinforcement learning approaches. We introduce \textit{Direct Preference Optimization (DPO)}, an algorithm that, while implicitly targeting the same objective as traditional RLHF techniques (maximization of reward subject to a KL-divergence regularization), offers greater implementation simplicity and ease of training. At a high level, the DPO update increases the ratio between the log-probabilities assigned to preferred versus non-preferred responses. It achieves this while integrating a dynamic, example-specific importance weighting mechanism, mitigating the risk of model collapse observed when naively maximizing the log-probability ratio. DPO, like prior approaches, depends on a theoretical framework for preference modeling—such as the Bradley-Terry model \cite{bradley1952rankanalysis}—which evaluates the agreement between reward signals and observed preference judgments. Contrasting with prior art, which employs the preference model to craft a loss function for reward model training (subsequently used to optimize the policy), DPO directly defines the preference loss over the policy by leveraging a variable substitution. As a result, with a collection of pairwise human preference data on model completions, DPO enables direct policy optimization via a straightforward binary cross-entropy loss, producing the optimal policy for an implicit reward inferred from preferences.

Our primary contribution is the proposal of Direct Preference Optimization (DPO), a simple algorithm for preference-based language model training that avoids reinforcement learning. Empirical results demonstrate that DPO matches or exceeds the performance of current techniques—including PPO-driven RLHF—on preference-driven tasks such as sentiment adjustment, text summarization, and dialogue, using models with up to 6B parameters.

\section{Related Work}

As self-supervised language models have grown in scale, they have demonstrated the capacity to solve some problems without explicit supervision, utilizing either zero-shot learning \citep{radford2019language} or few-shot prompting strategies \citep{gpt3,megatron,chowdhery2022palm}. Nonetheless, their effectiveness on specific downstream tasks and fidelity to user instructions can be considerably enhanced by fine-tuning with curated datasets composed of prompts and human-generated responses \citep{mishra-etal-2022-cross,sanh2022multitask,chung2022scaling,thoppilan2022lamda}. This fine-tuning process, known as `instruction-tuning,' not only expands the ability of large language models (LLMs) to adapt to unseen instructions but also improves their general utility \citep{chung2022scaling}. While instruction tuning has yielded strong outcomes, it is typically easier to gather comparative (i.e., relative) human evaluations of responses than to collect high-quality expert demonstrations. Building on this, recent research has pursued fine-tuning LLMs on collections of human preference data, leading to gains in diverse areas such as translation \citep{kreutzer-etal-2018-reliability}, summarization \citep{stiennon2022learning,ziegler2020finetuning}, narrative generation \citep{ziegler2020finetuning}, and instruction compliance \citep{ouyang2022training,ramamurthy2023is}. These approaches typically begin by training a neural reward model to align with preference datasets, often operationalized using a framework like the Bradley-Terry model \citep{bradley1952rankanalysis}. Subsequently, the language model is further refined to maximize the expected reward, employing reinforcement learning methods such as REINFORCE \citep{williams1992reinforce}, proximal policy optimization (PPO; \cite{schulman2017proximal}), or similar techniques \citep{ramamurthy2023is}. A related stream of research uses LLMs, previously refined through human-feedback-driven instruction following, to generate synthetic preference data focusing on qualities like safety or harmlessness, guided by minimal human supervision in the form of rubric-based annotations \citep{bai2022constitutional}. This direction synthesizes progress from two key research areas: the use of reinforcement learning for diverse language model objectives~\citep{Ranzato2015SequenceLT,paulus2018a,wu2018learning}, and advancements in learning directly from human preference data \citep{christiano2017deep,kupcsik2018learning}. Despite the intuitive value of leveraging human preference comparisons, using reinforcement learning to further fine-tune large language models introduces considerable practical challenges. The present work addresses these issues by presenting a theoretically grounded method for optimizing with respect to relative preferences without relying on reinforcement learning techniques.

Policy learning from preferences has been examined in both bandit and reinforcement learning paradigms outside the domain of language, giving rise to a variety of methodologies. In the contextual bandit scenario, when action choices are guided by preferences or rankings instead of explicit rewards, the problem is formulated as a contextual dueling bandit (CDB; \cite{yue2012karmed,dudik2015contextual}). Within CDBs, where absolute reward signals are unavailable, theoretical analysis often relies on the concept of a \textit{von Neumann winner}—a policy that is expected to outperform any other policy with probability at least 50\% \citep{dudik2015contextual}. Nonetheless, unlike human preference learning, where one learns from a static, offline collection of preference-labeled action pairs \citep{yan2022human}, CDBs involve obtaining preference feedback in an online fashion. Similarly, \textit{preference-based RL} (PbRL) involves learning through binary preferences determined by an \textit{unknown} 'scoring' function instead of receiving direct reward signals \citep{BusaFekete2014,ruiz2023dueling}. There is a suite of PbRL algorithms—some able to leverage off-policy preference datasets—but most approaches entail an explicit initial estimation of the unobserved scoring (reward) function, which is subsequently optimized \citep{jain2013learning,BusaFekete2014,christiano2017deep,sadigh2017active,kupcsik2018learning}. In contrast, we propose a unified policy optimization framework that directly aligns the policy with the observed preferences in a single learning phase.

\section{Preliminaries}\label{section:prelims}

Here we recapitulate the RLHF pipeline detailed in \citeauthor{ziegler2020finetuning} and later iterations \citep{stiennon2022learning, bai2022training, ouyang2022training}. This process is typically composed of three core steps: (1) supervised fine-tuning (SFT), (2) preference sampling and reward inference, and (3) reinforcement learning-based policy optimization.

\textbf{SFT}: The initial stage of RLHF involves refining a pre-trained language model by supervised learning on carefully curated datasets relevant to the target application (e.g., conversation, text summarization), yielding a model referred to as $\pi^\text{SFT}$.

\textbf{Reward Modelling Phase}: During this phase, the SFT-tuned model is prompted with inputs $x$, generating response pairs $(y_1, y_2)\sim \pi^\text{SFT}(y \mid x)$. Human annotators are then tasked with indicating a preference for one of the two outputs, denoted $y_w\succ y_l \mid x$, where $y_w$ is the preferred and $y_l$ the less preferred candidate from $(y_1, y_2)$. The mechanism generating these preferences is presumed to be an unobserved reward function $r^*(y, x)$. Several modeling strategies are utilized for the preference data, with the Bradley-Terry (BT) model \cite{bradley1952rankanalysis} serving as a standard choice (the framework can also accommodate broader models, such as Plackett-Luce, when multiple rankings are available \citep{plackett1975analysis, luce2012individual}). According to the BT model, the probability distribution $p^*$ representing human preferences is defined as follows:

Given a static collection of comparison data, denoted $\mathcal{D}=\bigl\{x^{(i)}, y_w^{(i)}, y_l^{(i)}\bigr\}_{i=1}^N$ and sampled according to $p^*$, we may define a reward function $r_{\phi}(x, y)$ parameterized by $\phi$ and estimate its parameters by maximizing the likelihood of the data. When recast as a binary classification task, this estimation is equivalent to minimizing the negative log-likelihood objective:

\begin{equation}\label{eq:reward_model}
    \mathcal{L}_R(r_{\phi}, \mathcal{D}) = -\mathbb{E}_{(x, y_w, y_l)\sim \mathcal{D}}\bigl[\log \sigma(r_{\phi}(x, y_w)- r_{\phi}(x, y_l))\bigr]
\end{equation}

where $\sigma$ represents the logistic sigmoid function. For large language models, the reward network $r_{\phi}(x, y)$ is typically initialized using the SFT policy $\pi^\text{SFT}(y \mid x)$, appending an additional linear head on the terminal transformer layer to output a scalar reward score~\cite{ziegler2020finetuning}. To promote reward stability, previous work has applied a normalization procedure so that $\mathbb{E}_{x,y\sim \mathcal{D}}\left[r_\phi(x, y)\right] = 0$ for every $x$.

\textbf{RL Fine-Tuning Phase}: In the subsequent RL stage, the policy is updated by utilizing the learned reward model for supervision. As described in previous research~\citep{jaques2017sequence, jaques2020human}, the objective is to maximize

\begin{equation}\label{eq:RL}
\max_{\pi_{\theta}}  \mathbb{E}_{x\sim \mathcal{D}, y\sim \pi_{\theta}(y \mid x)}\bigl[r_{\phi}(x, y)\bigr] - \beta\mathbb{D}_{\textrm{KL}}\bigl[\pi_{\theta}(y\mid x)\mid \mid \pi_\text{ref}(y\mid x)\bigr],
\end{equation}

where the hyperparameter $\beta$ adjusts the tolerance for deviation from a base reference policy, typically $\pi_\text{ref} = \pi^\text{SFT}$. Generally, the trainable policy $\pi_\theta$ is also initialized as $\pi^\text{SFT}$. Enforcing this constraint is critical to limit policy divergence from the data region where the reward model is reliable and to preserve output diversity, which helps avoid concentration on only a handful of outputs with highest reward. Due to the fact that outputs are sampled from discrete distributions in language modeling, the above optimization problem cannot be solved via gradient descent and necessitates the use of reinforcement learning methods. A prevailing strategy~\citep{ziegler2020finetuning, stiennon2022learning, bai2022training, ouyang2022training} is to define the augmented reward ${r(x, y) = r_{\phi}(x, y) -\beta (\log \pi_{\theta}(y\mid x) - \log \pi_\text{ref}(y\mid x))}$ and apply policy optimization techniques, such as PPO~\cite{schulman2017proximal}, to maximize this objective.

\section{Direct Preference Optimization}\label{sec:DPO}

Faced with the computational demands associated with running reinforcement learning at scale, particularly for fine-tuning language models, our approach seeks a more straightforward and efficient alternative for preference-based policy learning. Departing from the standard RLHF pipeline that separates reward model training from RL policy optimization, our method takes advantage of a specially chosen parameterization for the reward function that admits a closed-form solution for its optimal policy, eliminating the need for iterative RL. As elaborated in the following, the central insight is to employ an analytic correspondence between a reward and its optimal policy, facilitating a direct transformation of a loss defined on the reward model into a loss defined on the policy itself. This change-of-variables perspective circumvents the explicit fitting of a reward model, while allowing us to retain alignment with existing human preference models such as the Bradley-Terry framework. Consequently, the policy network concurrently models both the generation process and the (implicit) reward signal.

\textbf{DPO Objective Derivation.} To derive the DPO objective, we begin by considering the standard RL objective in Eq.~\ref{eq:RL} under a general reward $r$, in line with previous approaches. As established by prior work~\citep{peters2007reinforcement, peng2019advantage, korbak2022reinforcement, go2023aligning}, the optimal solution to the KL-regularized reward maximization problem in Eq.~\ref{eq:RL} admits the following solution:

\begin{equation}\label{eq:op_policy}
    \pi_r(y\mid x) = \frac{1}{Z(x)}\pi_\text{ref}(y\mid x)\exp\left(\frac{1}{\beta}r(x, y)\right),
\end{equation}

with $Z(x) =\sum_{y}\pi_\text{ref}(y\mid x)\exp\left(\frac{1}{\beta}r(x, y)\right)$ serving as the partition function. The detailed derivation is provided in Appendix~\ref{app:derivation1}. Even when employing the maximum likelihood estimator $r_\phi$ for the true reward function $r^*$, evaluating the partition function $Z(x)$ remains computationally intensive \citep{korbak2022reinforcement, go2023aligning}, which makes practical application of this form challenging. Nevertheless, we may algebraically rearrange Eq.~\ref{eq:op_policy} so as to represent the reward $r$ using the optimal policy $\pi_r$, the reference $\pi_\text{ref}$, and the unknown partition function $Z(\cdot)$. Taking logarithms on both sides and performing some algebraic manipulation leads to:

\begin{equation}\label{eq:main_eq}
    r(x,y) =\beta \log \frac{\pi_r(y\mid x)}{\pi_\text{ref}(y\mid x)} + \beta \log Z(x).
\end{equation}

This form allows us to re-express the ground-truth reward $r^*$ in terms of its optimal policy $\pi^*$. Notably, the Bradley-Terry model only requires the difference in reward values for two candidate outputs, i.e., ${p^*(y_1 \succ y_2 \mid x) = \sigma(r^*(x, y_1) - r^*(x, y_2))}$. Substituting Eq.~\ref{eq:main_eq} into the Bradley-Terry model in Eq.~\ref{eq:bradley-terry} results in the partition functions canceling out, allowing the human preference probability to be determined solely by $\pi^*$ and $\pi_\text{ref}$. Therefore, under the Bradley-Terry framework, the optimal RLHF policy $\pi^*$ obeys the following:

\begin{equation}\label{eq:objective}
    p^*(y_1\succ y_2 \mid x)=\frac{1}{1 + \exp\left(\beta \log \frac{\pi^*(y_2\mid x)}{\pi_\text{ref}(y_2\mid x)} - \beta \log \frac{\pi^*(y_1\mid x)}{\pi_\text{ref}(y_1\mid x)}\right)}
\end{equation}

Full details of the derivation can be found in Appendix~\ref{app:derivation2}. While Eq.~\ref{eq:objective} employs the Bradley-Terry formulation, similar results hold for broader preference models such as Plackett-Luce~\citep{plackett1975analysis, luce2012individual}, as outlined in Appendix~\ref{app:plackett_luce_models}.

Expressing the likelihood of human preference data directly in terms of the optimal policy rather than the reward, we can thus formulate a maximum likelihood criterion for a parameterized policy $\pi_\theta$. This policy-based loss, analogous to the reward modeling formulation (cf. Eq.~\ref{eq:reward_model}), is given by:

\begin{equation}\label{eq:optimum_model}
    \mathcal{L}_\text{DPO}(\pi_{\theta}; \pi_\text{ref}) = -\mathbb{E}_{(x, y_w, y_l)\sim \mathcal{D}}\left[\log \sigma \left(\beta \log \frac{\pi_{\theta}(y_w\mid x)}{\pi_\text{ref}(y_w\mid x)} - \beta \log \frac{\pi_{\theta}(y_l\mid x)}{\pi_\text

In this formulation, an implicit reward function is learned through an alternative parameterization, resulting in an optimal policy given directly by $\pi_\theta$. Additionally, this approach is mathematically equivalent to fitting a reparametrized Bradley-Terry model, thereby inheriting desirable theoretical characteristics, such as consistency under suitable conditions on the distribution of preference data \cite{bong2022generalized}. Further theoretical properties of DPO in relation to similar methodologies are explored in Section~\ref{sec:theory}.

\textbf{How does the DPO update function?} To gain mechanistic insight into DPO, it is instructive to inspect the gradient of its loss, $\mathcal{L}_\text{DPO}$. The derivative with respect to the parameters $\theta$ takes the following form:

\begin{multline*}\label{eq:gradient}
    \nabla_\theta \mathcal{L}_\text{DPO}(\pi_\theta;\pi_\text{ref}) = \\ -\beta\mathbb{E}_{(x, y_w, y_l) \sim \mathcal{D}} \bigg[\underbrace{\sigma(\hat{r}_\theta(x, y_l) - \hat{r}_\theta (x, y_w))}_\text{assigns more weight if the reward estimate is inaccurate}\bigg[\underbrace{\nabla_\theta\log \pi(y_w \mid x)}_\text{raise probability of $y_w$} - \underbrace{\nabla_\theta\log\pi(y_l \mid x)}_\text{lower probability of $y_l$}\bigg]\bigg],
\end{multline*}

Here, $\hat{r}_\theta(x, y) = \beta \log \frac{\pi_\theta(y \mid x)}{\pi_\text{ref}(y \mid x)}$ designates the reward signal implicitly characterized by $\pi_\theta$ relative to the reference model $\pi_\text{ref}$ (elaborated in Section~\ref{sec:theory}). In practical terms, the gradient of $\mathcal{L}_\text{DPO}$ encourages the model to increase the log-likelihood of preferred completions $y_w$ while reducing it for less-preferred responses $y_l$. Notably, the update for each example is scaled by the extent to which the current implicit reward model $\hat{r}_\theta$ overvalues dispreferred completions, proportional to $\beta$, which reflects the severity of reward misranking and is tied to the KL regularization parameter. Empirical evidence underscores the critical role of this weighting; omitting it (i.e., in a variant lacking this coefficient) can result in substantial degeneration of the model's language generation capabilities, as evidenced in Appendix Table~\ref{tab:unlikelihood_generations}.

\textbf{DPO overview.} The standard workflow for DPO involves the following steps: 1) For each prompt $x$, generate two completions $y_1, y_2 \sim \pi_\text{ref}(\cdot \mid x)$, use human judgments to assign preferences, and thereby build an offline dataset $\mathcal{D} = \{x^{(i)}, y_w^{(i)}, y_l^{(i)}\}_{i=1}^N$ of preference-labeled examples. 2) The language model $\pi_\theta$ is then trained by minimizing the DPO loss $\mathcal{L}_\text{DPO}$, using the specified $\pi_\text{ref}$, dataset $\mathcal{D}$, and chosen temperature parameter $\beta$. 

In applied settings, it is often more practical to utilize available public preference datasets rather than curate new ones through fresh sampling and annotation. Typically, these datasets originate from completions drawn from a supervised fine-tuned policy $\pi^\text{SFT}$, thus when possible, $\pi_\text{ref}$ is initialized to this $\pi^\text{SFT}$. If $\pi^\text{SFT}$ cannot be obtained, $\pi_\text{ref}$ is instead initialized by maximizing the likelihood on preferred completions, specifically via ${\pi_\text{ref} = \argmax_{\pi}\mathbb{E}_{x, y_w \sim \mathcal{D}}\left[\log \pi(y_w \mid x)\right]}$. This initialization scheme serves to reduce distributional discrepancy between the inaccessible true reference distribution and the chosen $\pi_\text{ref}$ for DPO training. Comprehensive details on experimental setup and hyperparameter selection are presented in Appendix~\ref{app:implementation}.

\section{Theoretical Analysis of DPO} In this section, we provide a deeper examination of DPO’s conceptual foundations, discuss supporting theoretical arguments, and explain how DPO addresses some deficiencies present in RLHF approaches (such as actor-critic algorithms like PPO~\cite{schulman2017proximal}).

\label{sec:theory}

\subsection{Your Language Model Is Secretly a Reward Model} DPO uniquely avoids constructing a separate reward model or engaging in explicit reinforcement learning; instead, it learns the target policy through a single maximum likelihood loss. Importantly, the objective in Eq.~\ref{eq:main_eq} matches a Bradley-Terry preference model, wherein the reward function is parameterized as $r^*(x, y) = \beta \log\frac{\pi^*_\theta(y \mid x)}{\pi_\text{ref}(y \mid x)}$. The optimization directly adapts $\pi_{\theta}$, paralleling reward model fitting in Eq.~\ref{eq:reward_model} under an appropriate reparameterization. This section develops the theoretical framework for this equivalence, demonstrates that it imposes no restrictions on the class of representable reward functions, and establishes the possibility of exactly recovering the optimal policy. To lay the groundwork, we introduce an equivalence relation on reward functions.

\begin{definition}
Two reward functions $r(x, y)$ and $r'(x, y)$ are considered equivalent if and only if there exists a function $f$ such that ${r(x, y)-r'(x, y) = f(x)}$.
\end{definition}

This definition naturally yields an equivalence relation, thereby dividing the collection of reward functions into distinct equivalence classes. From this, we obtain two fundamental lemmas:

\begin{lemma}\label{lemma:same_prefrence}
Within the context of the Plackett-Luce preference model, which includes the Bradley-Terry model as a special case, any two reward functions that belong to the same equivalence class result in identical preference distributions.
\end{lemma}

\begin{lemma}\label{lemma:same_policy}
    Any pair of reward functions within the same equivalence class yield the same optimal policy for the constrained reinforcement learning problem.
\end{lemma}

The arguments establishing these results are direct, and their details can be found in Appendix \ref{app:lemma1}. The first lemma highlights a well-documented issue of under-specification present in the Plackett-Luce model family \cite{plackett1975analysis}. As a consequence, it is standard practice to introduce additional identifiability conditions to ensure the MLE solutions to Eq. \ref{eq:reward_model} are meaningful \cite{bong2022generalized}. The second lemma indicates that the optimal policy is invariant to the choice of representative within an equivalence class. Therefore, for practical purposes, we only aim to identify any reward function from the correct equivalence class. We formalize this insight with the following theorem, whose proof appears in Appendix~\ref{app:thm1}:

\begin{theorem}\label{thm:main}
    Provided certain mild conditions, every reward equivalence class aligned with the Plackett-Luce (and Bradley-Terry) model assumptions may be described using the reparameterization ${r(x, y) = \beta \log \frac{\pi(y\mid x)}{\pi_\text{ref}(y\mid x)}}$, for some choice of model $\pi(y\mid x)$ and a predetermined reference model $\pi_\text{ref}(y \mid x)$.
\end{theorem}

\begin{sproof}
    Let $r(x, y)$ be an arbitrary reward function, inducing an optimal model $\pi_r(y \mid x)$ as characterized in Eq. \ref{eq:op_policy}. We aim to demonstrate that a representative from $r$'s equivalence class can be written in the reparameterized form stated above. Define the transformation $f$ as  
\begin{equation}
    f(r; \pi_\text{ref}, \beta)(x, y) = r(x, y) - \beta\log\sum_{y}\pi_\text{ref}(y\mid x)\exp\left(\frac{1}{\beta}r(x, y)\right)
\end{equation}
Here, $f$ rescales the reward by subtracting the log-partition function of $\pi_r$. Since this normalization term depends only on $x$, $f(r; \pi_\text{ref}, \beta)(x, y)$ remains within the equivalence class of $r(x, y)$. Substituting $r$ by the right side of Eq.~\ref{eq:main_eq}, which universally applies to any reward function, yields $f(r; \pi_\text{ref}, \beta)(x, y) = \beta \log \frac{\pi_r(y\mid x)}{\pi_\text{ref}(y\mid x)}$. Thus, $f$ maps $r$ to a member of its equivalence class with the desired form, ensuring that the reparameterization fully preserves the generality of the reward model.
\end{sproof}

An alternative interpretation of Theorem~\ref{thm:main} is that it precisely determines which reward function is chosen by the DPO reparameterization from within each equivalence class. Specifically, it selects the reward function that fulfills:

\begin{equation}\label{eq:lag_p}
     \sum_{y}\underbrace{\pi_\text{ref}(y\mid x)\exp\left(\frac{1}{\beta}r(x, y)\right)}_{=\pi(y\mid x)\text{, using Thm.~\ref{thm:main} reparam.}} = 1,
\end{equation}

meaning that $\pi(y\mid x)$ forms a proper probability distribution, with non-negative entries summing to unity. Observing Eq.~\ref{eq:op_policy}, it becomes evident that Eq.~\ref{eq:lag_p} corresponds to the partition function for the optimal policy as defined by the reward function $r(x, y)$. A fundamental contribution of the DPO algorithm is the ability to introduce particular constraints into the otherwise underdetermined Plackett-Luce (and more specifically, Bradley-Terry) models of preferences. By doing so, it retains the expressive capacity of the class of reward models, while also ensuring that the optimal policy in Eq.~\ref{eq:op_policy} remains computationally tractable for every prompt $x$.

\subsection{Instability of Actor-Critic Algorithms} Our framework additionally facilitates the identification of sources of instability inherent in conventional actor-critic methods applied in RLHF, including PPO. Concentrating on the RL fine-tuning component described in Section \ref{section:prelims}, we can relate this to the control as inference paradigm~\cite{levine2018reinforcement} within the context of the constrained RL formulation presented in \ref{eq:RL}. Given a parameterized policy $\pi_{\theta}(y\mid x)$, the objective is to minimize $\mathbb{D}_{\text{KL}}[\pi_{\theta}(y|x) \mid \mid \pi^*(y\mid x)]$, where $\pi^*$ denotes the optimal policy from Eq.~\ref{eq:optimum_model} as defined by the reward function $r_{\phi}(y, x)$. A straightforward algebraic manipulation yields the following optimization problem:

\begin{equation}\label{eq:AC}
    \max_{\pi_{\theta}}\mathbb{E}_{\pi_{\theta}(y\mid x)}\bigg[\underbrace{r_{\phi}(x, y) -\beta\log\sum_{y}\pi_\text{ref}(y\mid x)\exp\left(\frac{1}{\beta}r_{\phi}(x, y)\right)}_{f(r_{\phi}, \pi_\text{ref}, \beta)} - \underbrace{\beta\log\frac{\pi_{\theta}(y\mid x)}{\pi_\text{ref}(y\mid x)}}_{\text{KL}}\bigg]
\end{equation}

The objective discussed here is identical to that optimized in previous studies 
\citep{ziegler2020finetuning, stiennon2022learning, bai2022training, ouyang2022training}, which utilize the DPO-equivalent reward within the class $r_{\phi}$. Within this context, the normalization term in $f(r_{\phi}, \pi_\text{ref}, \beta)$ may be interpreted as representing the soft value function corresponding to the reference policy $\pi_\text{ref}$. Although the presence of this term does not influence the optimality of the solution, omitting it could lead to high variance in the policy gradient estimates and thus destabilize training. One approach to mitigate this normalization term involves learning a value function, though this can also present optimization challenges. Another strategy, employed in previous works, involves normalizing the rewards against a baseline provided by human completions, which effectively functions as a Monte Carlo estimate of the normalization constant based on a single sample. By contrast, the DPO reparameterization produces a reward formulation that eliminates the requirement for any external baselines.

\section{Experiments}
This section presents an empirical analysis of DPO's effectiveness in learning policies directly from preference data. We begin with a controlled text-generation experiment to investigate the efficiency of DPO in balancing reward maximization and KL-divergence minimization with the reference policy, benchmarking against prevalent preference optimization methods such as PPO. Subsequently, we assess DPO's performance on more complex RLHF tasks and with larger model scales, including tasks like summarization and dialogue. Our results indicate that DPO, with minimal hyperparameter adjustment, consistently matches or surpasses robust baselines such as RLHF with PPO and approaches using the selection of the highest reward trajectory out of $N$ samples under a learned reward model. Prior to detailing these outcomes, we outline the experimental methodology; supplementary information can be found in Appendix~\ref{app:exp_details}.

\textbf{Tasks.} We investigate three distinct open-ended text generation tasks within our experimental framework. In each case, algorithms are trained to optimize a policy using a preference dataset $\mathcal{D}=\bigl\{x^{(i)}, y_w^{(i)}, y_l^{(i)}\bigr\}_{i=1}^N$. The first task, \textbf{controlled sentiment generation}, considers $x$ as the beginning segment of a movie review selected from the IMDb dataset \cite{maas-EtAl:2011:ACL-HLT2011}; here, the policy is tasked with producing a continuation $y$ that expresses positive sentiment. To enable controlled assessment, we construct preference pairs automatically with a pre-trained sentiment classifier, selecting those pairs for which $p(\text{positive}\mid x,y_w)>p(\text{positive}\mid x,y_l)$. For supervised fine-tuning (SFT), GPT-2-large is further trained to convergence on the IMDb training reviews (see App~\ref{app:sentiment_details} for implementation specifics). The second task, \textbf{summarization}, treats $x$ as a Reddit forum post requiring the policy to generate a summary $y$ capturing key information. Consistent with earlier research, we make use of the Reddit TL;DR summarization dataset \citep{volske-etal-2017-tl} as well as human-annotated preferences from \citeauthor{stiennon2022learning}. The supervised baseline leverages an SFT model, trained on human-generated summaries of forum posts\footnote{\url{https://huggingface.co/CarperAI/openai_summarize_tldr_sft}}, and further optimized via RLHF using the TRLX framework \citep{leandro_von_werra_2023_7790115}. The preference labels correspond to model completions from a comparably trained, albeit separate, SFT model, as described by \citeauthor{stiennon2022learning}. For the final task, \textbf{single-turn dialogue}, $x$ is an arbitrary user prompt ranging from scientific inquiries to personal advice. The objective is to produce a reply $y$ that is both useful and engaging. Data comes from the Anthropic Helpful and Harmless dialogue corpus \citep{bai2022training}, which includes 170,000 dialogues, each ending with two assistant responses and an associated human preference indicator. Since no ready-made SFT model exists for this corpus, we construct one by fine-tuning a general-purpose language model solely on the preferred completions.

\textbf{Evaluation.} We utilize two distinct evaluation frameworks in our experiments. For the controlled sentiment generation scenario, where the ground-truth reward function (a sentiment classifier) is available, we assess each algorithm’s ability to optimize the constrained reward maximization objective by measuring its performance frontier with respect to both achieved reward and KL-divergence from the reference policy. This analysis is enabled by direct access to the actual reward function. In contrast, real-world applications lack explicit access to the true reward function, necessitating alternative evaluation criteria. Thus, we report each algorithm’s \textit{win rate} compared to a baseline policy, using GPT-4 as an approximate surrogate for human assessment of summary quality and response helpfulness in summarization and single-turn dialogue, respectively. In summarization, test set reference summaries serve as baselines, while for dialogue tasks, the test dataset’s preferred responses are employed. Despite literature indicating that language models are often superior automated evaluators relative to traditional metrics \citep{Chen2023ExploringTU}, we further validate the reliability of GPT-4 judgments by conducting a human study, as discussed in Sec.~\ref{sec:human-judgments}. Our results indicate a strong correlation between GPT-4 and human evaluations, with agreement rates between human annotators and GPT-4 generally matching or exceeding those among human annotators themselves.

\textbf{Methods.} Alongside DPO, we systematically assess various established methods for training language models to align with human preferences. As a straightforward baseline, we implement zero-shot prompting with \textbf{GPT-J} \citep{gpt-j} for summarization and employ 2-shot prompting using \textbf{Pythia-2.8B} \citep{biderman2023pythia} for the dialogue task. Additionally, we examine the \textbf{SFT} model, as well as \textbf{Preferred-FT}—a model trained via supervised learning on the chosen completion $y_w$, generated either by the SFT model (for sentiment and summarization) or by a standard language model (in dialogue tasks). Another considered approach is the pseudo-supervised \textbf{Unlikelihood}~\citep{welleck2019neural} method, which trains policies to simultaneously maximize the likelihood of $y_w$ and suppress the likelihood of $y_l$, modulated by a tunable coefficient $\alpha\in[0,1]$ on the unlikelihood objective. We also include \textbf{PPO} \citep{schulman2017proximal}, trained using rewards inferred from preference data, and \textbf{PPO-GT}, an oracle variant that utilizes the ground truth reward in the controlled sentiment setting. In sentiment experiments, two versions of PPO-GT are utilized: a standard implementation \cite{leandro_von_werra_2023_7790115} and a modified variant featuring reward normalization and enhanced hyperparameter tuning to bolster performance—modifications also applied when executing standard PPO with learned rewards. Lastly, the \textbf{Best of $N$} baseline is investigated, which samples $N$ responses from either the SFT model (or Preferred-FT in the case of dialogue) and selects the top response according to a reward model trained on preference data. While this method can deliver strong results by separating reward modeling from PPO optimization, its practicality is limited since inference for each query requires generating $N$ completions, rendering it computationally burdensome for all but small $N$.

\subsection{How well can DPO optimize the RLHF objective?}

In standard RLHF methodologies, the KL-constrained reward maximization objective ensures that while maximizing reward, the learned policy does not diverge excessively from the reference policy. Accordingly, meaningful comparison of different algorithms necessitates evaluating both the obtained reward and the associated KL divergence; notably, increasing reward at the expense of disproportionately high KL is not inherently advantageous. As depicted in Figure~\ref{fig:frontier-tldr-main}, we present the reward-KL tradeoff achieved by several methods in the sentiment classification scenario. For each method, we conduct multiple training experiments, varying a key hyperparameter controlling the policy's conservativeness—specifically, using target KL values $\in\{3,6,9,12\}$ for PPO, $\beta \in \{0.05,0.1,1,5\}$, and $\alpha\in\{0.05,0.1,0.5,1\}$ for unlikelihood training, with preferred-FT evaluated across different random seeds. This comprehensive sweep encompasses 22 total runs. At every interval of 100 training steps until convergence, we assess the resultant policies using a collection of held-out test prompts, calculating both the mean reward per the true reward signal and the mean sequence-level KL\footnote{Specifically, the aggregate of the KL divergences computed at each timestep.} relative to the reference policy, $\text{KL}\left(\pi\mid \mid \pi_\text{ref}\right)$. The experimental outcomes indicate that DPO attains a frontier that is significantly more efficient, reaching the highest rewards while maintaining low KL divergence. This finding is particularly remarkable for several reasons. Firstly, despite DPO and PPO targeting identical optimization objectives, DPO attains superior efficiency; DPO’s reward-versus-KL tradeoff unambiguously outperforms that of PPO. Secondly, DPO surpasses PPO in the reward-KL landscape \emph{even when PPO is granted access to true reward labels} (PPO-GT).

\subsection{Can DPO scale to real preference datasets?}
\label{sec:dpo-real-datasets}
We proceed to investigate the effectiveness of DPO fine-tuning on tasks involving summarization and single-turn dialogue. In the context of summarization, it is well-established that metrics like ROUGE often display weak alignment with human judgment~\citep{stiennon2022learning}. Previous research has demonstrated that preference-based fine-tuning of language models with PPO leads to more desirable summaries. In our assessment, we compare several methods by generating completions on the TL;DR summarization dataset’s test split and calculating their average win rates against the dataset’s reference summaries. For each method, completions are sampled across a range of temperatures from 0.0 to 1.0, and win rate results are illustrated in Figure~\ref{fig:frontier-tldr-main} (right). DPO, PPO, and Preferred-FT all start from an identical GPT-J SFT checkpoint\footnote{\url{https://huggingface.co/CarperAI/openai_summarize_tldr_sft}}. Our experiments show that at temperature 0.0, DPO achieves an approximate win rate of 61\%, surpassing PPO's optimal performance of roughly 57\% at the same temperature. Additionally, DPO attains a greater maximum win rate than the best-of-$N$ approach. It is important to point out that DPO's $\beta$ hyperparameter was not subject to extensive tuning, suggesting that these outcomes may understate its capabilities. Furthermore, we observe that DPO maintains more consistent performance across varying sampling temperatures than PPO, whose results decline to match the base GPT-J model at higher temperatures. In contrast, Preferred-FT does not yield marked improvements over SFT. Section~\ref{sec:human-judgments} further presents direct human evaluation, where DPO outputs sampled at temperature 0.25 are preferred 58\% of the time over PPO outputs sampled at the same temperature.

For the single-turn dialogue setting, we conduct assessments using the segment of the Anthropic HH test set \citep{bai2022training} involving just one exchange between a human and an assistant. The evaluation with GPT-4 leverages the test set’s preferred completions as ground truth to calculate win rates for each approach. Due to the lack of an established SFT benchmark for this scenario, our workflow begins with a pre-trained Pythia-2.8B model; we then apply Preferred-FT, training a reference model on the selected completions so that the outputs stay close to the model’s expected distribution, before proceeding with DPO training. We include as baselines the best result among 128 Preferred-FT completions (after noting that the Best of $N$ baseline reaches its performance ceiling at $N=128$; refer to Appendix Figure~\ref{fig:best-of-n}) and a 2-shot prompt configuration for the original Pythia-2.8B model. Our experiments show that DPO equals or exceeds the performance of these alternatives, contingent on optimal temperature selection for each. Additionally, we test an RLHF model trained using PPO on the Anthropic HH dataset \footnote{\url{https://huggingface.co/reciprocate/ppo_hh_pythia-6B}} and sourced from a widely recognized implementation \footnote{\url{https://github.com/CarperAI/trlx/tree/main/examples/hh}}, yet we are unable to achieve results surpassing those of the plain Pythia-2.8B model regardless of prompt design or sampling temperature. Given our TL;DR findings and recognizing that both methods are tuned to the same reward function, we treat the Best of 128 performance as an approximate stand-in for PPO outcomes. In sum, DPO stands out as the only method that not only remains computationally efficient but also outperforms the preferred completions on the Anthropic HH dataset, while achieving parity or superiority compared to the much more resource-intensive Best of 128 approach. As depicted in Figure~\ref{fig:dialogue-main}, DPO rapidly attains its peak performance.

\subsection{Generalization to a new input distribution}

To further explore how PPO and DPO handle distributional changes, we assess both algorithms—using the policies trained in our Reddit TL;DR summarization study—on a novel input: news articles from the test portion of the CNN/DailyMail dataset \citep{nallapati-etal-2016-abstractive}. For sampling, we employ the optimal temperatures identified for TL;DR (0 and 0.25). The corresponding outcomes are summarized in Table~\ref{tab:ood}. We utilize the same GPT-4 (C) prompt as in the TL;DR evaluation, substituting “forum post” with “news article,” to determine GPT-4's win rate when benchmarked against the reference summaries from the dataset. Within this distinct domain, DPO continues to deliver notably stronger results compared to PPO. These findings suggest that DPO’s generalization capabilities are on par with, if not superior to, those of PPO—even though DPO does not leverage the additional unlabelled Reddit TL;DR prompts available to PPO.

\subsection{Validating GPT-4 judgments with human judgments}
\label{sec:human-judgments}
To assess the trustworthiness of GPT-4 as an evaluator, we run a human subject study referencing the TL;DR summarization experiment and utilizing two variant prompts for GPT-4. The \textbf{GPT-4 (S)} (simple) prompt inquires which summary captures the essential content more effectively. In contrast, the \textbf{GPT-4 (C)} (concise) prompt adds a criterion for conciseness, motivated by observations that GPT-4, under the simple prompt, tends to select longer, sometimes redundant summaries more often than human raters. (See Appendix~\ref{app:prompts} for the precise phrasing.) Our evaluation involves three model-policy comparisons: the top-performing (DPO at temperature 0.25), lowest-performing (PPO at temperature 1.0), and an intermediate baseline (SFT at temperature 0.25), each matched against PPO with greedy sampling at its best temperature—intended to reflect a spectrum of generation qualities. Analysis reveals that GPT-4’s agreement with human preferences, across both prompts, aligns closely with the agreement rate observed among human raters, supporting its utility as an automated proxy for human assessment (note: due to limited human resources, multiple annotations are obtained only for DPO and PPO-1). Results also show that the \textbf{GPT-4 (C)} prompt produces win rates most aligned with human judgments, so we rely on it for the main analyses in Section~\ref{sec:dpo-real-datasets}. Further methodological details, such as the web interface for raters and the roster of participating volunteers, can be found in Appendix~\ref{app:human-study}.

\section{Discussion}
Preference-based learning represents an effective and scalable approach for developing advanced, value-aligned language models. In this work, we presented DPO, a straightforward training strategy enabling preference-driven language model optimization without resorting to reinforcement learning. Unlike conventional approaches that reformulate preference learning as a reinforcement learning problem to apply standard RL methods, DPO leverages a correspondence between the policy of a language model and its underlying reward function. This approach allows for direct alignment with human preferences via a standard cross-entropy objective, bypassing both reinforcement learning machinery and any compromise in model generality. Remarkably, DPO achieves competitive, and sometimes superior, performance relative to established RLHF techniques such as those built on PPO, requiring minimal hyperparameter adjustment. Consequently, DPO significantly lowers the entry point for training language models guided by human preferences.

\textbf{Limitations \& Future Work.} Several open questions remain in light of our findings. For instance, how does generalization to out-of-distribution data compare between the DPO policy and models trained on explicit reward signals? Preliminary evidence indicates that DPO-based policies can generalize on par with those from PPO training, though more systematic analysis is warranted. An additional avenue to explore is whether utilizing self-labeling mechanisms with the DPO policy can yield similar benefits for prompts lacking explicit annotations. Furthermore, understanding the dynamics of reward over-optimization within direct preference optimization—such as interpreting the mild performance dip seen in Figure~\ref{fig:dialogue-main}-right—merits further investigation. Our evaluation considered models up to 6B parameters, yet there is substantial promise in applying DPO at the scale of the largest current language models. With respect to assessment, we observe that GPT-4-derived win rates are sensitive to prompt formulation, highlighting the need for refined methods of obtaining reliable judgments from automated evaluators. Lastly, DPO has substantial potential beyond its application to language models; it could be adapted to optimize generative models across various data modalities.